\section{Setting and evaluation protocol}
\label{sec:setting}
\paragraph{Problem.}
A history of $M$ models has been scored on $N$ benchmarks; the score matrix $X\in\mathbb{R}^{M\times N}$ may have missing cells (mask $\Omega$).
A new model can be run on $k$ benchmarks $S$ chosen from those it can be run on, and the evaluator wants its scores on the other benchmarks $T$.
We call $k^\ast(\epsilon)$ the smallest $k$ at which the error on $T$ falls to a tolerance $\epsilon$.
Scores are standardised per benchmark with the mean and standard deviation of the history rows only, so that errors are in units of a benchmark's standard deviation and benchmarks on different scales (accuracy, win rate, nDCG) are comparable.

\paragraph{Protocol.}
The unit of replication is the held-out \emph{model}.
The models are split into history and test models by repeated 5-fold partitions; for every split, the standardisation, the factor model and its rank, all hyper-parameters (chosen by cross-validation inside the history), the selection order, and the conformal calibration use history rows only.
A test model is scored on a benchmark only if it has a score there: on sparse matrices the candidates for $S$ are the test model's own observed benchmarks and $T$ is its remaining observed benchmarks, so no score is ever imputed for evaluation.
Confidence intervals are 95\% percentile bootstrap intervals over held-out models, after averaging over repeated splits (and over random draws for the random selector).
We report, at each $k$: RMSE in standardised units over $T$; the pairwise accuracy of the question ``which history models does the new model beat'' on each target, averaged over targets; top-10 F1 (does the predicted ranking recover the history models in the top ten); and Spearman correlation across a model's targets.
Fitting on all rows, with the test models in the history, makes the error look smaller by $0.04$--$0.10$ in MSE ($10$--$23\%$ in RMSE) on MTEB, with 95\% intervals that exclude zero in every cell (Appendix~\ref{app:protocol}).

\paragraph{Matrices.}
Table~\ref{tab:t1} (Appendix~\ref{app:data}) lists the matrices.
The primary data are the MTEB text-embedding leaderboard results \citep{muennighoff2022mteb} at a pinned commit (CC0-1.0): the ``models'' are embedding models and the ``benchmarks'' are embedding tasks, not generative LLMs.
Two settings are fixed from the missingness pattern alone, before any error was computed: a complete block of 103 models $\times$ 202 tasks, found by repeatedly dropping the lowest-coverage row or column (A1), and the 126 models complete on the 30 most-covered tasks (A2); and the sparse matrix of 507 models $\times$ 598 tasks with 22\% of cells observed, with the availability constraint above (B).
As a second kind of matrix we use generative-LLM leaderboards (Open LLM Leaderboard v2 \citep{fourrier2024open}, HELM Lite and HELM MMLU \citep{liang2022helm}, LiveBench \citep{white2024livebench}) with $N$ between 6 and 57; for these, held-out folds are whole model families or organisations.
Synthetic matrices with known loadings serve to check the theory (Section~\ref{sec:theory}).

\paragraph{Comparisons.}
The confirmatory comparisons were written down before the confirmatory runs (budgets $k\in\{3,8,20\}$): the greedy selector against random choice and against clustering, pivoted-QR, leverage-score and centrality selectors; the factor posterior against ridge, kNN and softImpute \citep{mazumder2010spectral}; cost-aware against cost-ignorant selection; conformal against model-based stopping.
They are Holm-corrected \citep{holm1979simple} as one family per experiment and matrix; everything else is labelled exploratory.
Two comparisons were added after a pilot on a separate seed had shown that the first criterion degrades at large $k$ (Section~\ref{sec:experiments}); they are marked $\dagger$.
